% !TEX root = 0000.tex
% Anexo B
% Contenido tomado de main_2.tex (líneas 3579-3684) sin cambios de texto.

% El texto está en inglés: reglas de corte de palabra del inglés
\begin{hyphenrules}{english}

\chapter{Anexo: Energy Estimates and Convergence Bound for Hyperbolic PDEs}\chaptermark{Energy Estimates}\label{app:derivations}

% ---------------------------------------------------------------------
\section{Energy Estimate for a Klein--Gordon PINN}\label{app:energy}

We prove Theorem~\ref{thm:energy-estimate}. Multiply the error equation \eqref{eq:error-pde},
$\partial_t^2e-\Delta e+m^2e=r$, by $\partial_te$ and integrate over $\Omega$. We analyze each term
separately.

\paragraph{Term with $\partial_t^2e$.}
Since $\partial_t\big((\partial_te)^2\big)=2\,\partial_te\,\partial_t^2e$,
\[
  \int_\Omega\partial_t^2e\,\partial_te\,d\mathbf{x}=\frac12\frac{\,d}{\,d t}\int_\Omega(\partial_te)^2\,d\mathbf{x}.
\]

\paragraph{Term with $-\Delta e$ (spatial integration by parts).}
We apply Green's first identity:
\[
  -\int_\Omega\Delta e\,\partial_te\,d\mathbf{x}=\int_\Omega\nabla e\cdot\nabla(\partial_te)\,d\mathbf{x}
  -\int_{\partial\Omega}\frac{\partial e}{\partial\mathbf{n}}\,\partial_te\,d S,
\]
where $\partial e/\partial\mathbf{n}=\nabla e\cdot\mathbf{n}$ is the outward normal derivative. In addition,
$\nabla e\cdot\nabla(\partial_te)=\frac12\partial_t\abs{\nabla e}^2$; therefore,
\begin{equation}
  -\int_\Omega\Delta e\,\partial_te\,d\mathbf{x}=\frac12\frac{\,d}{\,d t}\int_\Omega\abs{\nabla e}^2\,d\mathbf{x}
  -\int_{\partial\Omega}\frac{\partial e}{\partial\mathbf{n}}\,\partial_te\,d S.
  \label{eq:app-lap}
\end{equation}

\paragraph{Term with $m^2e$.}
\[
  \int_\Omega m^2e\,\partial_te\,d\mathbf{x}=\frac12\frac{\,d}{\,d t}\int_\Omega m^2e^2\,d\mathbf{x}.
\]

\paragraph{Assembly and energy identity.}
Adding the three terms and using the energy $E(t)$ of Equation~\eqref{eq:error-energy}, the integrated equation is
\begin{equation}
  \frac{\,d E}{\,d t}=\int_\Omega r\,\partial_te\,d\mathbf{x}+\int_{\partial\Omega}\frac{\partial e}{\partial\mathbf{n}}\,\partial_te\,d S.
  \label{eq:app-identity}
\end{equation}
With $r=0$ and periodic or homogeneous boundary conditions, this is the conservation of energy
Equation~\eqref{eq:kg-energy}.

\paragraph{Bounds for the source and boundary terms.}
For the interior term we use Young's inequality ($ab\le\frac12a^2+\frac12b^2$):
\[
  \int_\Omega r\,\partial_te\,d\mathbf{x}\le\frac12\norm{r}^2_{L^2(\Omega)}+\frac12\norm{\partial_te}^2_{L^2(\Omega)}
  \le\frac12\norm{r}^2_{L^2(\Omega)}+E(t),
\]
since $\frac12\norm{\partial_te}^2_{L^2}\le E(t)$. For the boundary integral we use that
$\partial_te|_{\partial\Omega}=\partial_tB$ and apply Young's inequality again:
\[
  \int_{\partial\Omega}\frac{\partial e}{\partial\mathbf{n}}\,\partial_tB\,d S
  \le\frac12\Big\|\frac{\partial e}{\partial\mathbf{n}}\Big\|^2_{L^2(\partial\Omega)}+\frac12\norm{\partial_tB}^2_{L^2(\partial\Omega)}.
\]
Substituting in Equation~\eqref{eq:app-identity} we obtain
\begin{equation}
  \frac{\,d E}{\,d t}\le E(t)+\frac12\Big(\norm{r}^2_{L^2(\Omega)}+\norm{\partial_tB}^2_{L^2(\partial\Omega)}
  +\Big\|\frac{\partial e}{\partial\mathbf{n}}\Big\|^2_{L^2(\partial\Omega)}\Big)=:E(t)+\frac12\rho(t).
  \label{eq:app-differential}
\end{equation}

\paragraph{Gronwall's inequality.}
Multiplying by $e^{-t}$, $\frac{\,d}{\,d t}\big(e^{-t}E(t)\big)\le\frac12e^{-t}\rho(t)\le\frac12\rho(t)$, and
integrating from $0$ to $t$,
\[
  E(t)\le\Big(E(0)+\frac12\int_0^t\rho(s)\,d s\Big)e^{t},
\]
which is Equation~\eqref{eq:energy-bound}. (A cruder use of Young's inequality gives the factor $e^{2t}$ with a constant $C$, which is also
a valid, but weaker, bound.) The initial energy is
\[
  E(0)=\frac12\int_\Omega\Big(\abs{I_1}^2+\abs{\nabla I_0}^2+m^2\abs{I_0}^2\Big)\,d\mathbf{x}
  =\frac12\Big(\norm{I_1}^2_{L^2}+\norm{\nabla I_0}^2_{L^2}+m^2\norm{I_0}^2_{L^2}\Big).
\]
With this, we obtain a bound of the energy of the error in terms of the $L^2$ norms of the residuals, of the
gradient of the initial position residual, and of the boundary terms. In practice, for the convergence of
the PINN in the energy norm, the loss must control $\norm{r}^2_{L^2}$, $\norm{I_1}^2_{L^2}$,
$\norm{I_0}^2_{L^2}$ and $\norm{\nabla I_0}^2_{L^2}$ (plus the boundary terms for non-periodic boundaries),
which guarantees $E(t)\to0$ and, therefore, convergence in $H^1$; in one spatial dimension, by the Sobolev
embedding $H^1\subset C^{0,1/2}$, the convergence is also uniform.

% ---------------------------------------------------------------------
\section{From the Dirac Equation to the Klein--Gordon Equation}\label{app:dirac}

Since each component of the Dirac spinor individually satisfies the Klein--Gordon equation
\cite{bjorken1964}, the study of the dynamics of wave packets and relativistic dispersion can be done
validly in the scalar formalism of Klein--Gordon. Let the Dirac equation (natural units, metric signature
$(+,-,-,-)$) be
\begin{equation}
  (i\gamma^{\mu}\partial_\mu-m)\psi=0.
  \label{eq:dirac}
\end{equation}
We apply the operator $(i\gamma^{\nu}\partial_\nu+m)$:
\begin{align*}
  0&=(i\gamma^{\nu}\partial_\nu+m)(i\gamma^{\mu}\partial_\mu-m)\psi\\
   &=\underbrace{-\gamma^{\nu}\gamma^{\mu}\partial_\nu\partial_\mu\psi}_{\text{kinetic term}}
     \underbrace{-\,i m\gamma^{\nu}\partial_\nu\psi+i m\gamma^{\mu}\partial_\mu\psi}_{\text{cancel}}
     \underbrace{-\,m^2\psi}_{\text{mass term}}.
\end{align*}
Since $\partial_\nu\partial_\mu$ is symmetric, only the symmetric part of $\gamma^{\nu}\gamma^{\mu}$ survives,
$\frac12\{\gamma^{\nu},\gamma^{\mu}\}=g^{\nu\mu}$, and $g^{\nu\mu}\partial_\nu\partial_\mu=\partial^{\mu}\partial_\mu$.
So we are left with
\[
  (\partial^{\mu}\partial_\mu+m^2)\psi=0\iff(\partial_t^2-\nabla^2+m^2)\psi=0,
\]
and we arrive at the Klein--Gordon equation for each component of $\psi$.

\end{hyphenrules}
